\section{Combining the odd and mixed normalizations}
\label{sec:dyadic-metabelian-intersection}

The arithmetic commutators in degrees eight and ten have trivial beta
multiplier. They therefore let us combine the infinite odd normalization
with the finite leading-term normalization.

\begin{corollary}\label{cor:dyadic-metabelian-deep-gate}
Put
\[
 E=C_\infty\cap D_{10},\qquad J_E=E\cap H.
\]
Then $E$ is a proper closed normal subgroup of $D_{10}$ and of
infinite index there, and
\[
 G=EH,\qquad B\simeq E/\overline{[G,E]J_E}.
\]
Every pro-$2$ Magnus word on $E$ starts in degree at least eleven,
and its dyadic beta multiplier is identically one.
The kernel of restriction from $H^1(N,\mathbf F_2)^G$ to
$H^1(E,\mathbf F_2)$ is exactly
\[
 \left(\bigoplus_{q\text{ odd prime}}\mathbf F_2\psi_q\right)
                   \oplus\mathbf F_2\eta_3.
\]
\end{corollary}

\begin{proof}
We first verify the factorization without assuming any joint-image
statement. Since $C_\infty\subseteq D_7$, the coordinate $m_8$
is defined there. The arithmetic word $w_8$ of the eighth
normalization is a commutator of multiplier-one elements. The
beta multiplier is a homomorphism on $M=\ker\lambda$, so
$B_{w_8}=1$. Hence $w_8\in C_\infty\cap H$, and its $m_8$
value is a unit. Correcting by its powers gives
\[
 C_\infty=(C_\infty\cap D_8)(C_\infty\cap H).
\]
Since $m_9=0$ on $C_\infty$, the first factor lies in $D_9$.
The arithmetic word $w_{10}$ of the tenth normalization is also
a commutator of multiplier-one elements, so belongs to
$C_\infty\cap J_9$ and has an odd $m_{10}$ value. Correcting
by its powers gives
\[
 C_\infty\cap D_8=E(C_\infty\cap D_8\cap H).
\]
Together with $G=C_\infty H$, these equations prove $G=EH$.
Normality and the relative quotient follow as before.

To prove strictness, take the word $v_5\in C_4\cap H$ in
\eqref{eq:uniform-odd-arithmetic-witness}. It has
$m_{11}(v_5)\in2^{11}\mathbf Z_2^{\times}$.
First correct its eighth coordinate by a power of $w_8$.
The result lies in $C_4\cap D_8\subseteq D_9$.
Then correct its tenth coordinate by a power of $w_{10}$.
The result $v$ belongs to $D_{10}\cap H$.
Both correcting words have beta multiplier one, so they leave
$\tau_{11}$, and therefore the nonzero eleventh depth-one
coefficient, unchanged. Thus $v\notin E$.
In fact $D_{10}\subseteq C_4$, so the uniform bound and the
corrected arithmetic word give the exact images
\[
 m_{11}(D_{10})=m_{11}(D_{10}\cap H)=2^{11}\mathbf Z_2.
\]
This additive coordinate vanishes on $E$, proving infinite index. The final assertions follow
from $E\subseteq D_{10}$ and its definition.

For the last kernel assertion, every invariant character $\chi$ on
$N$ kills both correcting words. Indeed,
\[
 w_8=[\sigma,h\sigma h^{-1}]=[\sigma,k_8],\qquad
 k_8=(h\sigma h^{-1})\sigma^{-t}\in D\subseteq N;
\]
the rational Kummer coordinates of $k_8$ cancel by their exact
weight-one law. Invariance under $\sigma\in G$ gives
$\chi(w_8)=0$. The tenth construction gives
$w_{10}=[\sigma,k_{10}]$ for $k_{10}\in D$ and hence
$\chi(w_{10})=0$ as well. The two correction factorizations in
this proof now imply that $\chi$ vanishes on $E$ if and only if
it vanishes on $C_\infty$. Proposition~\ref{prop:uniform-odd-character-kernels}
gives the asserted kernel.
\end{proof}
